%% kapa: exam=1 round=36 year=2025 period=2 form=A subject=appraisal_law subject_ko=감정평가관계법규 seq=21 number=21
%% source: Q-Net 감정평가사 기출문제 (cst00309) artlSeq=5231525 file=raw/1008/5231525/2025년 제36회 감정평가사 1차시험 2교시 문제지 원본.pdf
%% conversion: kordoc 4.24.51 -> tools/gen1.py (auto); review=proofread-sonnet; flags=none
\begin{question}{21}
국유재산법령상 행정재산의 사용허가에 관한 설명으로 옳지 않은 것은?
\begin{choices}
  \item 보존용재산을 사용허가하는 경우에 재산의 유지ㆍ보존을 위하여 관리비가 특히 필요할 때에는 사용료에서 그 관리비 상당액을 뺀 나머지 금액을 징수할 수 있다.
  \item 행정재산이 토지의 용도 등을 고려할 때 해당 재산에 인접한 토지의 소유자를 지명하여 경쟁에 부칠 필요가 있는 경우에는 지명경쟁의 방법으로 사용허가를 받을 자를 결정할 수 있다.
  \item 중앙관서의 장은 공공용 재산을 사용허가하는 경우 그의 용도나 목적에 장애가 되지 아니하는 범위에서만 사용허가를 할 수 있다.
  \item 행정재산으로 할 목적으로 기부를 받은 재산의 상속인에게 사용허가를 하는 경우 사용허가기간은 5년 이내로 한다.
  \item 중앙관서의 장은 행정재산의 사용허가를 받은 자가 해당 재산의 보존을 게을리한 경우에는 그 허가를 취소하거나 철회할 수 있다.
\end{choices}
\end{question}
\answer{4}
